\section{Results}
\label{sec:results}

\subsection{Aggregate coordination}
Figure~\ref{fig:profiles} shows the cold-day aggregate load. The uncoordinated
profile exhibits the classic double peak --- a morning and an evening heating
surge --- whereas the ideal ADMM coordination flattens it to a nearly constant
$\approx 362.3$~kW. The $\varrho = 0.8$ trace, in which most agents are imputed,
recovers much of the flattening but leaves residual peaks where imputation error
accumulates.

\begin{figure}[t]
\centering
\includegraphics[width=\columnwidth]{fig_aggregate_profiles.pdf}
\caption{Aggregate cold-day load: uncoordinated double-peak versus the flat
ADMM-ideal profile ($\approx 362.3$~kW) and the $\varrho = 0.8$ imputed case with
residual peaks. Here $\varrho$ is the non-responsive fraction.}
\label{fig:profiles}
\end{figure}

Table~\ref{tab:agg} quantifies the aggregate key performance indicators. Ideal
coordination cuts the peak from 401.6 to 362.3~kW, a 9.8\% reduction, and drives
the peak-to-average ratio (PAR) from 1.109 down to 1.000 --- essentially a flat
load. A PAR of 1.000 is the practical floor: it means the coordinated aggregate
is constant to within a tenth of a percent across the entire day, so the
flattening objective in~\eqref{eq:sharing} has been met almost exactly. Energy
cost falls slightly, from 909.5 to 895.4~CAD --- a 1.55\% reduction --- because
the daily-energy constraint conserves total heating; the saving comes from
shifting load out of the evening-peak TOU window (0.158~CAD/kWh) toward off-peak
hours (0.078~CAD/kWh) rather than from using less energy. The modest size of the
cost reduction is itself a useful result, and an honest one: the objective
in~\eqref{eq:sharing} flattens the aggregate (it penalizes departure from a
\emph{constant} level), it does not chase the TOU price signal. A purely flat
profile still places substantial load in the priced evening window, so the bill
moves only as a side effect of peak-shaving. With energy conserved, demand
response of this kind buys \emph{capacity} relief (a lower peak) far more than it
buys energy savings, and the case for it must rest on the network benefit rather
than the bill.

\subsubsection{Comfort preservation}
The flattening is achieved without sacrificing thermal comfort. The mean absolute
heating deviation of the ideal coordinated schedule from the thermostat baseline
is 0.264~kW per home-step --- small against the $\approx 5$~kW heating level the
homes operate at. This is comfort preserved \emph{by construction}: the
projection~\eqref{eq:proj} keeps every step inside the $\pm 50\%$ band of its
baseline and conserves each home's daily heating energy, so the coordinator can
only re-time heat, never withhold it. The result is a feeder-level flat profile
that is, at the level of each individual home, only a gentle nudge of when the
thermostat draws power.

\subsubsection{A plateau-and-ramp transition}
Swept finely (Table~\ref{tab:agg}), the degradation of coordination with the
non-responsive fraction is not a smooth decay but a \emph{plateau followed by a
ramp}. From $\varrho = 0$ to $\varrho \approx 0.70$ the coordinated peak stays
clamped at the ideal optimum (362.3--362.6~kW, PAR within $0.1\%$ of unity);
beyond a knee near $\varrho = 0.72$ it rises approximately linearly toward the
uncoordinated peak, reaching 368.5, 375.0, 389.6, and finally 401.9~kW at
$\varrho = 0.75, 0.8, 0.9$, and $1.0$. This shape is the signature of a
\emph{flexibility-budget} mechanism, and it has two causes. First, an imputed
home is pinned to its forecast --- which reproduces its \emph{natural}, peaky
heating shape, not a flattened one --- so it leaves the coordination pool;
degradation is therefore driven mainly by the \emph{loss of controllable
flexibility}, not by raw forecast error (which stays small, RMSE $\approx 0.3$~kW,
last column). Second, the remaining responsive homes flatten the total only up to
their collective $\pm 50\%$ deferral budget: as long as that budget exceeds the
peakiness injected by the pinned population, the peak is held exactly at the
optimum (the plateau); once it is exhausted --- around $\varrho \approx 0.72$,
where only $\sim 20$ of 74 homes remain controllable --- the residual peak passes
through, and it does so linearly. The transition is thus sharp at its
\emph{onset} (a binding-constraint knee) but smooth thereafter (a linear ramp);
the earlier impression of a sudden step was an artifact of coarse sampling
between $\varrho = 0.7$ and $1.0$.

\begin{table}[t]
\centering
\caption{Aggregate KPIs by non-responsive fraction $\varrho$ (selected points;
the full 12-point sweep appears in Figure~\ref{fig:peakpar}). Iter.: ADMM
iterations to converge.}
\label{tab:agg}
\begin{tabular}{@{}lrrrrr@{}}
\toprule
Scenario & Peak (kW) & PAR & Cost (CAD) & RMSE (kW) & Iter. \\
\midrule
Uncoordinated      & 401.6 & 1.1090 & 909.5 & ---   & ---  \\
Coordinated ideal  & 362.3 & 1.0000 & 895.4 & ---   & 8    \\
$\varrho{=}0.4$  & 362.4 & 1.0003 & 895.4 & 0.305 & 18   \\
$\varrho{=}0.55$ & 362.4 & 1.0004 & 895.5 & 0.302 & 27   \\
$\varrho{=}0.6$  & 362.4 & 1.0004 & 895.5 & 0.304 & 28   \\
$\varrho{=}0.65$ & 362.5 & 1.0006 & 895.5 & 0.303 & 36   \\
$\varrho{=}0.7$  & 362.6 & 1.0008 & 895.5 & 0.301 & 65   \\
$\varrho{=}0.75$ & 368.5 & 1.0171 & 895.8 & 0.305 & 300  \\
$\varrho{=}0.8$  & 375.0 & 1.0350 & 896.6 & 0.304 & 300  \\
$\varrho{=}0.9$  & 389.6 & 1.0753 & 902.6 & 0.303 & 300  \\
$\varrho{=}1.0$  & 401.9 & 1.1093 & 909.5 & 0.304 & 300  \\
\bottomrule
\end{tabular}
\end{table}

\subsection{Imputation accuracy and robustness}
The LightGBM imputer is a competent forecaster of an unseen home's heating. Under
5-fold leave-agents-out cross-validation it achieves an RMSE of 0.304~kW and an
MAE of 0.255~kW, against a heating standard deviation of 0.736~kW. At roughly
41\% of the natural variability, the imputer is clearly better than predicting
the mean and is accurate enough that the coordinated aggregate barely moves for
non-responsive fractions below the flexibility knee, as seen in the per-scenario
RMSE column of Table~\ref{tab:agg}. Notably the per-scenario imputation RMSE is
nearly flat in $\varrho$ ($\approx 0.30$~kW) even as the aggregate degrades sharply ---
direct evidence that the damage past the knee comes from the loss of controllable
agents, not from worsening forecasts.

Figure~\ref{fig:peakpar} plots peak and PAR against $\varrho$ over the full
12-point sweep. Both remain clamped at the optimum through $\varrho \approx 0.70$
and then rise approximately linearly to the uncoordinated values at $\varrho = 1$
--- a plateau-and-ramp, not a gradual decay, with the ramp onset (knee) marking
exhaustion of the responsive deferral budget.

\begin{figure}[t]
\centering
\includegraphics[width=\columnwidth]{fig_peak_par_vs_rho.pdf}
\caption{Aggregate peak and PAR versus the non-responsive fraction $\varrho$. The
horizontal axis labelled ``rho'' in the plotted image denotes this fraction
$\varrho$, not the ADMM penalty $\rho_{\mathrm{ADMM}}$.}
\label{fig:peakpar}
\end{figure}

A secondary finding concerns tractability rather than outcome quality. The ADMM
iteration count grows steeply with $\varrho$: 8 iterations for the ideal case,
then 18, 27, 28, and 36 through $\varrho = 0.65$, rising to 65 at $\varrho = 0.70$
and then hitting the 300-iteration cap at $\varrho = 0.75$ and beyond (all runs
below the cap converged). Communication failure thus degrades coordination
\emph{tractability} --- the effort to reach consensus --- well before it degrades
the coordinated \emph{outcome}, an early-warning signal an operator could
monitor.

The shape of this curve is informative, and it shares its knee with the peak and
the network results. Through $\varrho = 0.70$ the iteration count stays bounded
($\le 65$); at $\varrho = 0.75$ it more than quadruples to the 300 cap even though
the peak there has only just begun to move (368.5~kW). The knee of the iteration
curve therefore \emph{precedes} the visible outcome degradation: the solver feels
the flexibility budget run out --- the sharing residual becomes hard to drive to
zero as the imputed aggregate diverges from the one that physically forms --- one
sweep step before the peak and the transformer loading break. Iteration count is
thus a cheap, purely software-side proxy for the harder-to-observe
physical-feasibility margin, visible before any meter or power-flow check, because
both respond to the same underlying quantity: the mismatch injected by
imputation.

\subsection{Imputation method comparison and the value of imputation}
The results so far use a single imputer (LightGBM). Two questions remain: how
much does imputation \emph{itself} buy, and how much does the \emph{choice} of
method matter? To answer them fairly we adopt a stricter \emph{realized-load}
model that separates belief from reality: a silent home physically draws its
\emph{true} (uncontrolled, peaky) heating, while the coordinator sees only an
\emph{estimate} and flattens the responsive homes against it; the feeder then
experiences the realized aggregate --- responsive coordinated schedules plus the
true silent load plus background. ``No imputation'' is the coordinator flattening
the responsive subset \emph{blind} to the silent homes. We compare four learned
estimators (LightGBM, random forest, ridge regression, $k$-nearest neighbours),
a naive ``mean-level'' baseline (a flat profile at each silent home's known
nameplate level --- correct energy, no shape), and the no-imputation case.

Table~\ref{tab:imputers} reports both layers of metric. Intrinsically (5-fold
leave-agents-out), ridge regression is the most accurate (RMSE 0.280~kW,
$R^2 = 0.70$), slightly ahead of LightGBM (0.304~kW) and $k$-NN (0.308~kW) and the
random forest (0.342~kW); with this small population and four smooth features the
linear model is hard to beat. Notably the naive mean-level baseline is close to
the learned methods on RMSE (0.312~kW) --- because each home's heating is
dominated by its mean level, which that baseline gets exactly right --- so
intrinsic RMSE alone badly under-rates the importance of getting the \emph{timing}
right.

\begin{table}[t]
\centering
\caption{Imputation methods: intrinsic accuracy (leave-agents-out CV) versus
realized feeder impact at $\varrho = 0.6$ (mean over 200 random silent subsets).
Lower realized loading and violation probability are better.}
\label{tab:imputers}
\begin{tabular}{@{}lrrrrr@{}}
\toprule
Method & RMSE & $R^2$ & Peak & Load. & $P_{\mathrm{viol}}$ \\
 & (kW) & & (kW) & (\%) & \\
\midrule
Ridge              & 0.280 & 0.70 & 366.9 &  99.8 & 0.17 \\
$k$-NN             & 0.308 & 0.67 & 368.1 & 100.1 & 0.61 \\
LightGBM           & 0.304 & 0.67 & 370.4 & 100.8 & 0.92 \\
Random forest      & 0.342 & 0.62 & 369.4 & 100.5 & 0.86 \\
Mean-level (naive) & 0.312 & 0.65 & 376.6 & 102.5 & 1.00 \\
No imputation      & ---   & ---  & 389.0 & 106.0 & 1.00 \\
\bottomrule
\end{tabular}
\end{table}

The realized-impact columns and Figure~\ref{fig:comparison} tell the operational
story, and here the violation probability $P_{\mathrm{viol}}$ is the sharpest
discriminator. Without imputation the realized peak climbs to 389.0~kW (106\%
loading) at $\varrho = 0.6$ and the transformer violates in \emph{every} draw
($P_{\mathrm{viol}} = 1.00$), because the outcome depends entirely on \emph{which}
uncontrolled homes fall silent. Any learned imputer holds the realized peak near
367--370~kW ($\approx 100$--101\%): imputation buys roughly a 19~kW reduction
(389 to 370~kW) and, more importantly, drops the violation probability from 1.00
(none) to as low as 0.17 (ridge). The naive mean-level baseline recovers part of
the peak (376.6~kW) but still violates in every draw --- direct evidence that
reconstructing the \emph{shape}, not merely the energy, is what lets the
responsive homes counter the silent peaks. Among the learned methods
$P_{\mathrm{viol}}$ now discriminates cleanly and tracks intrinsic forecast
quality: ridge 0.17, $k$-NN 0.61, random forest 0.86, and LightGBM 0.92, with the
realized-loading ranking none ($106\%$) $>$ naive ($102.5\%$) $>$ the ML cluster
($\approx 100$--101\%) $>$ ideal ($98.5\%$). The practical lesson is therefore
twofold: \emph{using} an imputer matters far more than \emph{which} one --- one
versus none is the gap that dominates --- and a simple ridge model is the best,
cheapest choice on data of this size, while LightGBM remains a robust default that
scales to the larger populations a utility would face.

\begin{figure}[t]
\centering
\includegraphics[width=\columnwidth]{fig_imputer_comparison.pdf}
\caption{Realized worst transformer loading versus the non-responsive fraction for
each imputation method under the realized-load model. No imputation (red) diverges
early; the learned methods cluster near the 100\% limit; the naive mean-level
baseline sits between them. All converge to the uncoordinated 109.6\% once every
home is silent.}
\label{fig:comparison}
\end{figure}

\subsection{Network validation}
The closed-loop contribution is summarized in Figure~\ref{fig:network} and
Table~\ref{tab:net}. The uncoordinated aggregate loads the 400~kVA transformer to
109.6\% --- a thermal violation --- while the worst minimum voltage sits at a
healthy 0.962~p.u. Ideal coordination \emph{clears} the thermal overload: maximum
transformer loading falls to 98.5\% (no violation) and the worst minimum voltage
stays healthy at 0.966~p.u. This is the core result: flattening the aggregate is
not merely a statistical improvement, it restores feeder feasibility.

\begin{figure}[t]
\centering
\includegraphics[width=\columnwidth]{fig_network_violation_vs_rho.pdf}
\caption{Worst minimum bus voltage (left axis) and worst maximum transformer
loading (right axis) versus the non-responsive fraction $\varrho$ --- the headline
feasibility figure. Loading stays cleared below 100\% through $\varrho = 0.70$,
crosses the limit at $\varrho = 0.75$, and ramps linearly to the uncoordinated
109.6\% at $\varrho = 1$. The image's ``rho'' axis is the non-responsive
fraction $\varrho$.}
\label{fig:network}
\end{figure}

\begin{table}[t]
\centering
\caption{Network feasibility on the LV feeder (worst of day), selected fractions.
Loading is the 400~kVA transformer's worst-of-day loading.}
\label{tab:net}
\begin{tabular}{@{}lrrrl@{}}
\toprule
Scenario & Peak (kW) & $V_{\min}$ (p.u.) & Loading (\%) & Violation \\
\midrule
Uncoordinated      & 401.6 & 0.962 & 109.6 & yes \\
Coordinated ideal  & 362.3 & 0.966 &  98.5 & no \\
$\varrho{=}0.4$  & 362.4 & 0.966 &  98.5 & no \\
$\varrho{=}0.55$ & 362.4 & 0.966 &  98.5 & no \\
$\varrho{=}0.7$  & 362.6 & 0.966 &  98.6 & no \\
$\varrho{=}0.75$ & 368.5 & 0.966 & 100.2 & yes \\
$\varrho{=}0.8$  & 375.0 & 0.965 & 102.1 & yes \\
$\varrho{=}0.9$  & 389.6 & 0.964 & 106.2 & yes \\
$\varrho{=}1.0$  & 401.9 & 0.962 & 109.6 & yes \\
\bottomrule
\end{tabular}
\end{table}

The cleared state is resilient to communication failure through $\varrho = 0.70$,
where the transformer holds at 98.6\% loading and does not violate. The
deterministic crossing falls between $\varrho = 0.70$ (98.6\%, cleared) and
$\varrho = 0.75$ (100.2\%, violated); beyond it the loading ramps linearly ---
102.1\%, 106.2\% --- back to the uncoordinated 109.6\% once every home is
imputed at $\varrho = 1$. The story is therefore
$109.6\% \rightarrow 98.5\% \rightarrow 109.6\%$: coordination clears the
overload, the clearance survives substantial non-responsiveness, and then the
protection unwinds linearly as the controllable population vanishes. This pins the
critical non-responsive fraction --- beyond which imputation can no longer protect
the transformer --- to a narrow band around $\varrho = 0.72$.

\paragraph{Voltage stays healthy.} Across every scenario the worst minimum LV bus
voltage stays in a comfortable $0.962$--$0.966$~p.u.\ band and never approaches
the 0.95~p.u.\ limit, so the feeder is voltage-comfortable throughout and the
\emph{transformer thermal loading} is the unambiguous binding constraint --- a
cleaner story than a benchmark that is stressed on voltage to begin with. The
100\% transformer rating gives an unambiguous binary verdict, and coordination is
what moves the feeder across it from infeasible to feasible.

\subsection{Forecast uncertainty quantification}
The network verdict above rests on a single realization per non-responsive
fraction. To test whether that point estimate understates risk, we run the
Monte-Carlo analysis of Section~\ref{sec:methodology} (200 draws per fraction over
which agents fall silent and Gaussian forecast residuals of standard deviation
0.304~kW, the imputer's cross-validated RMSE), mapping each draw's daily peak to
the worst-of-day transformer loading.

Figure~\ref{fig:uncertainty} and Table~\ref{tab:uncertainty} report the result
(mean with [P5, P95] band). The probability of a thermal violation stays at 0
through $\varrho = 0.65$, ticks to 0.04 at $\varrho = 0.70$, jumps to 0.96 at
$\varrho = 0.75$, and is 1.00 for $\varrho \ge 0.80$. The single-realization
analysis reported $\varrho = 0.70$ as comfortably compliant (98.6\% loading), but
once the uncertainty in which agents fail and in the forecast is accounted for,
the transformer still clears at $\varrho = 0.70$ (only 4\% of draws breach) yet
exceeds its limit in 96\% of draws by $\varrho = 0.75$: the deterministic estimate
understated how abruptly the risk rises. The [P5, P95] band is tight on the
plateau --- forecast errors average out across a large responsive population ---
and widens at the knee ($[98.6, 99.95]\%$ at $\varrho = 0.70$, $[100.1, 102.0]\%$
at $\varrho = 0.75$) before riding the ramp upward. This reframes the operator
takeaway: the critical transition is not a single deterministic fraction but a
probabilistic risk threshold centered near $\varrho = 0.72$, where the breach
probability swings from near zero to near certainty across a single sweep step.

\begin{figure}[t]
\centering
\includegraphics[width=\columnwidth]{fig_uncertainty_loading.pdf}
\caption{Monte-Carlo forecast-uncertainty sweep (200 draws per fraction). Worst
maximum transformer loading mean with its P5--P95 band versus the non-responsive
fraction $\varrho$ (left axis, with the 100\% thermal limit), and the probability
of a transformer violation as bars (right axis). The image's ``rho'' axis is the
non-responsive fraction $\varrho$.}
\label{fig:uncertainty}
\end{figure}

\begin{table}[t]
\centering
\caption{Monte-Carlo transformer loading and violation probability by fraction
(200 draws each; mean [P5, P95]).}
\label{tab:uncertainty}
\begin{tabular}{@{}lll@{}}
\toprule
$\varrho$ & Loading (\%) [P5, P95] & $P(\text{violation})$ \\
\midrule
0.40 & 98.5 [98.5, 98.5]    & 0.00 \\
0.55 & 98.5 [98.5, 98.5]    & 0.00 \\
0.70 & 99.1 [98.6, 99.95]   & 0.04 \\
0.75 & 101.0 [100.1, 102.0] & 0.96 \\
0.80 & 102.5 [101.6, 103.4] & 1.00 \\
1.00 & 110.0 [109.2, 110.9] & 1.00 \\
\bottomrule
\end{tabular}
\end{table}

\subsection{Discussion}
Three observations tie the results together. First, the aggregate, tractability,
and feasibility curves all share the same knee near $\varrho = 0.72$: the
iteration count, the peak, and the worst transformer loading all stay benign on
the plateau through $\varrho \approx 0.70$, break at the same knee, and ramp
together thereafter. This suggests that a single observable quantity --- the
divergence between the imputed and the realized aggregate, governed by the
responsive flexibility budget --- drives all three, and that an operator
monitoring ADMM convergence effort has an early proxy for impending feeder
violation (the iteration knee leads the loading knee by one sweep step). Second,
the network result is not a restatement of the aggregate result: a 9.8\% peak cut
is a statistical improvement, but \emph{clearing a 109.6\% transformer overload to
98.5\%} is an operational state change from infeasible to feasible, and the two
need not coincide on a different feeder or load allocation. Third, the unwinding
past the knee is graceful rather than catastrophic --- the loading ramps linearly
to 109.6\% and the power flow always converges, never collapsing --- which
indicates that the imputer keeps the system in a recoverable regime even when most
agents are silent. Together these support the central claim that coordination
quality should be measured by feeder feasibility, and that imputation extends the
range of communication failure over which that feasibility holds --- here, up to a
critical non-responsive fraction near $\varrho = 0.72$.
